\documentclass[12pt]{beamer}
\usepackage{beamerthemeCopenhagen}
\usecolortheme[rgb={0.18,0.3,0.3}]{structure}
%\usepackage[T1]{fontenc}
\usepackage[utf8]{inputenc}
%\usepackage[latin1]{inputenc}
\usepackage[brazil]{babel}
\usepackage{amsmath}
\usepackage{amsfonts}
\usepackage{amsthm}
\usepackage{pgf,tikz}
\usepackage{bbm}

\usetikzlibrary{arrows}

\newcommand{\caixa}{$\Box$}
\newtheorem{defi}{Definição}
\newtheorem{ex}{Exemplo}
\newtheorem{teo}{Teorema}
\newtheorem{cor}{Corolário}
\newtheorem{lema}{Lema}
\newtheorem{dist}{Distribuição}
\newtheorem{emv}{Máxima verossimilhança}
\newtheorem{mo1}{Primeiro momento}
\newtheorem{mo2}{Segundo momento}
\newtheorem{mc2}{Segundo momento central}
\newtheorem{prova}{Demonstração}
\newtheorem{BiNeg}{Definição}

\institute[IC-UFAL]{Universidade Federal de Alagoas\\
Instituto de Computação}
\title{Comparação de estimadores pelo método Monte Carlo das distribuições: $\mathcal{G}_{\rm I}^ {0}$, Uniforme, Binomial negativa, Gamma e Exponencial}
\author{Antônio Marcos Larangeiras\and Alisson Nascimento\and
  Paulo\and Bruno Normande\and Juliana Leal}
\date{{\tiny Maceió-AL, Dezembro de 2013}}
\begin{document}
\begin{frame}
\maketitle
\begin{picture}(0,0)(0,-50)
\includegraphics[scale=0.11]{IC.png}
\end{picture}
\begin{picture}(0,0)(-250,-40)
\includegraphics[scale=0.13]{ufal.png}
\end{picture}
\end{frame} 

%=========================
\section{Introdução}
%=========================


%================================================================================
\begin{frame}
\frametitle{Introdução}
\begin{itemize}
\item Calcular o estimador por máxima verossimilhança;

\bigskip

\item Calcular estimador pelo primeiro momento;

\bigskip

\item Calcular estimador pelo segundo momento;

\bigskip

\item Utilizar o Método Monte Carlo força bruta.

\end{itemize}
\end{frame}
%===============================================================================


%========================
\section{Distribuições}
%========================


%==============================================================================================
\subsection{Distribuição $\mathcal{G}_{\rm I}^{0}$}

\begin{frame}
\frametitle{Distribui\c c\~ ao $\mathcal{G}_{\rm I}^{0}$}

Denotaremos esta distribuição por $Z\sim\mathcal{G}_{\rm I}^{0}(z; \alpha,1,1)$, com $-\alpha, z>0$. Sua densidade é

\begin{displaymath}
f_{Z}=(z;\alpha,1,1)=\dfrac{\Gamma(L-\alpha)}{\Gamma(-\alpha)(1+z)^{1-\alpha}}, \mbox{com} -\alpha \ \mbox{e}\ z>0.
\end{displaymath}


\end{frame}

%========================================================================================================================================================
\begin{frame}
\frametitle{Distribui\c c\~ ao $\mathcal{G}_{\rm I}^{0}$}

A esperança matemática da distribuição $\mathcal{G}_{\rm I}^{0}$:

\begin{equation}
\label{eq:Gi0_esp}
E(Z)=\dfrac{\gamma}{\alpha-1}
\end{equation}\pause

A variância da distribuição $\mathcal{G}_{\rm I}^{0}$:

\begin{equation}
\label{eq:Gi0_var}
Var(Z)=\dfrac{1}{(\alpha-1)^2}+\dfrac{1}{(\alpha-2)(\alpha-1)^2}+\dfrac{1}{\alpha-1}
\end{equation}

\end{frame}

%====================
\subsubsection{Estimadores}
%====================
\begin{frame}
\frametitle{Estimadores da Distribuição $\mathcal{G}_{\rm I}^{0}$}
Pela Máxima Veressimilhança:
\begin{equation}
\label{eq:Gi0_MV}
\hat{\alpha}=\dfrac{n}{\displaystyle\sum_{i=1}^{n}\ln(1+z_i)}
\end{equation} \pause
Pelo momento central de ordem 2:
\begin{eqnarray}
\label{eq:Gi0_m_p2}
\hat{\alpha}_2=\overline{Z}
\end{eqnarray}
\end{frame}

\begin{frame}
\frametitle{Estimadores da Distribuição $\mathcal{G}_{\rm I}^{0}$}
Pelo segundo momento:
\begin{equation}
\label{eq:Gi0_m_p3}
\hat{\alpha}_3 = \frac{\frac{3}{k}\displaystyle\sum_{i=1}^{k}Z_{i}^{2}+2+\sqrt{\frac{1}{k^{2}}\displaystyle\sum_{i=1}^{k}Z_{i}^{4}+4}}{\frac{2}{k}\displaystyle\sum_{i=1}^{k}Z_{i}^{2}}
\end{equation}

\end{frame}


%===================
\subsubsection{Gráficos}
%=========================================================================================================================================================
\begin{frame}
\frametitle{Gráficos da distribuição $\mathcal{G}_{\rm I}^{0}$}

\end{frame}
%=========================================================================================================================================================




\subsection{Distribuição Uniforme}
%========================================================================================================================================================
\begin{frame}
\frametitle{Distribuição Uniforme $U(0, \theta)$}
\begin{dist}[$U(0, \theta)$]
Denotaremos esta distribuição por $X\sim U(x; 0,\theta)$, com $\theta>0$. Sua densidade é

{\small $$f(x) =
\left \{
\begin{array}{cc}
\frac{1}{ \theta } & 0 \leqslant x \leqslant \theta\\
0 & \textnormal{caso contrário.} \\
\end{array}
\right.$$ }

\bigskip\bigskip

{\scriptsize $E(X)=\dfrac{\theta}{2}$\hspace{1cm} e \hspace{1cm}$Var(X)=\dfrac{\theta^{2}}{12}$ }
\end{dist}
\end{frame}
%========================================================================================================================================================

%====================
\subsubsection{Estimadores}
%====================

\begin{frame}
\frametitle{Estimadores da distribuiç\~ao $U(0, \theta)$}
\begin{emv}[$U(0, \theta)$]
$$\widehat{\theta}_{1}={\dfrac{1}{\theta ^ n }}{\mathbb{1}_{\max{\{X_1,\dots,X_n\}\leqslant\theta}}},
\textnormal{portanto } \widehat{\theta} ={\max{\{X_1,\dots,X_n\}\leqslant\theta}} $$
\end{emv}

\begin{mo1}[$U(0, \theta)$]
$$\widehat{\theta}_{2}=2\overline{X}$$
\end{mo1} 

\end{frame}

\begin{frame}
\begin{mo2}[$U(0, \theta)$]
$$\widehat{\theta}_{3} = \sqrt{{{{3} \over {n}} {\sum_{i=1}^{n} X_{i}^{2}}}} $$
\end{mo2}
\end{frame}
%========================================================================================================================================================
%===================
\subsubsection{Gr\' aficos}
%===================
%=========================================================================================================================================================
\begin{frame}
\frametitle{Gráficos da distribuiç\~ao $U(0, \theta)$}

\end{frame}
%=========================================================================================================================================================

\subsection{Distribuição Binomial Negativa}
\begin{frame}
\frametitle{Distribuição Binomial Negativa}
Considere a situação de observar um evento dicotômico $Y_{i}$ independentes e identicamente distribuídos segundo uma lei de Bernoulli de probabilidade $p$. Suponha que se registre $Y$, o número de ensaios até obter exatamente $k$ sucessos. 

\end{frame}

\begin{frame}
\frametitle{Distribuição Binomial Negativa}

Seja uma variável aleatória que fornece o numero de ensaios até o k-ésimo sucesso. 
Assim, $Y$ tem uma distribuição binomial negativa com parâmetro $p \in (0,1)$, 
se sua função de probabilidade é dada por:

\begin{displaymath}
P_{r}(Y=y)=\left\{ \begin{array}{ll}
\binom{y-1}{k-1} \cdot p^{k} \cdot (1-p)^{y-k} & \textrm{se  y=k,k+1, \ldots.}\\
0 & \textrm{, caso contrário}
\end{array} \right.
\end{displaymath}


\end{frame}

%========================================================================================================================================================
\begin{frame}
\frametitle{Distribui\c c\~ ao $BN(p,k)$}

Usualmente sua função de probabilidade denota-se: $Y \sim BN(p,k)$.\pause

A esperança matemática da distribuição binomial negativa:

\begin{equation}
\label{eq:esp}
E(X) = {\frac{k}{p}}
\end{equation}\pause

A variância da distribuição binomail negativa:

\begin{equation}
\label{eq:var}
Var(X) = {\frac{k(1-p)}{p^{2}}}
\end{equation}

\end{frame}

%====================
\subsubsection{Estimadores}
%====================
\begin{frame}
\frametitle{Estimadores da Distribuição $BN(p,k)$}
Pela Máxima Veressimilhança:
\begin{equation}
\label{eq:BN_MV}
{\hat{p}} = {\frac{n k}{\sum_{i=1}^{n}{Y_{i}}}}
\end{equation} \pause
Pelo momento central de ordem 2:
\begin{eqnarray}
\label{eq:BN_m_p2}
{\hat{p}_{2}} = {\frac{-k + \sqrt{k^{2}+4k{Var(X_{i})}}}{2{Var(X_{i})}}}
\end{eqnarray}
\end{frame}

\begin{frame}
\frametitle{Estimadores da Distribuição $BN(p,k)$}
Pelo segundo momento:
\begin{equation}
\label{eq:BN_m_p3}
{\widehat{p}_{3}} = {\frac{{-k} + {\sqrt{k^{2} + 4 E(X_{i}^{2})(k^{2}+k)}}}{2E(X_{i}^{2})}}
\end{equation}

\end{frame}


%===================
\subsubsection{Gráficos}
%=========================================================================================================================================================
\begin{frame}
\frametitle{Gráficos da distribuição $BN(p,k)$}

\end{frame}
%=========================================================================================================================================================


\subsection{Distribuição Gamma}

% =====================================================================================
 \begin{frame}
   \frametitle{Distribuição Gamma}

   A Distribuição Gamma é caracterizada por dois valores, denominados
   \textit{shape} ($k$) e \textit{scale} ($\theta$).

   \begin{dist}[Função Probabilidade Densidade]
     $$ f(x;k,\theta) = \frac{x^{k-1}e^{-\frac{x}{\theta}}}{\theta^k\Gamma(x)} $$ 

     \begin{center}
       para   

       $ x > 0 $ e $ k,\theta > 0 $ 
     \end{center}
   \end{dist}
 \end{frame}
% ======================================================================================

% ====================
\subsubsection{Estimadores}
% ====================


\begin{frame}
  \frametitle{Estimadores da distribuição Gamma}

  Como foi definido que o valor de $\theta$ seria sempre igual a 1,
  então basta estimar os valores de $k$.

  \begin{emv}[Gamma]
    
  \end{emv}
\end{frame}
% ===================================================================================
\begin{frame}
  \frametitle{Estimadores da distribuição Gamma}
  \begin{mo1}[Gamma]
    \begin{align*}
      E[X] &= k\theta \\
      \hat{k} &= \frac{1}{n} \sum_{i=1}^nx_i
    \end{align*}
  \end{mo1}
  
  \begin{mc2}[Gamma]
    \begin{align*}
      Var[X] &= k\theta^2 \\
      \hat{k} &= Var(x)
    \end{align*}
  \end{mc2}
\end{frame}
% ========================================================================================================================================================

% ===================
\subsubsection{Gráficos}
% ===================
% =========================================================================================================================================================
\begin{frame}
  \frametitle{Gráficos da distribuição $\mathcal{G}_{\rm I}^{0}$}

\end{frame}
% =========================================================================================================================================================


%===================
\section{Resultados}
%===================

%================================================================================
\begin{frame}
\frametitle{Resultados}

\end{frame}
%================================================================================

%=====================
\section{Conclusões}
%=====================

%================================================================================
\begin{frame}
\frametitle{Conclusões}

\end{frame}
%===============================================================================

\section{Referências Bibliográficas}

\begin{frame}
\frametitle{Referências Bibliográficas}


\bibliographystyle{plain}

\bibliography{Larangeiras}


\end{frame}
\end{document}